\documentclass[11pt]{article}
\usepackage[margin=1in]{geometry}
\usepackage{amsmath,amssymb,amsthm,mathtools,bm}

%% ── shortcuts ──────────────────────────────────────────────
\newcommand{\X}{\mathbf{X}}
\newcommand{\y}{\mathbf{y}}
\newcommand{\bbR}{\mathbb{R}}
\newcommand{\one}{\mathbf{1}}
\newcommand{\tr}{\operatorname{tr}}
\newcommand{\Cov}{\operatorname{Cov}}
\newcommand{\Var}{\operatorname{Var}}
\newcommand{\E}{\mathbb{E}}
\newcommand{\bSig}{\boldsymbol{\Sigma}}
\newcommand{\balpha}{\boldsymbol{\alpha}}
\newcommand{\beps}{\boldsymbol{\epsilon}}
\newcommand{\bbeta}{\boldsymbol{\beta}}
\newcommand{\bx}{\mathbf{x}}
\newcommand{\hbSig}{\widehat{\boldsymbol{\Sigma}}}

\theoremstyle{plain}
\newtheorem{theorem}{Theorem}
\newtheorem{lemma}[theorem]{Lemma}
\newtheorem{proposition}[theorem]{Proposition}
\theoremstyle{definition}
\newtheorem{definition}[theorem]{Definition}
\theoremstyle{remark}
\newtheorem{remark}[theorem]{Remark}

\title{Variance of the Centered Oracle LDSC Estimator}
\author{}
\date{}

\begin{document}
\maketitle

%% ═══════════════════════════════════════════════════════════
\section{Setup and Notation}
%% ═══════════════════════════════════════════════════════════

Let $\X\in\bbR^{n\times p}$ be the genotype matrix whose rows
$\bx_1,\dots,\bx_n$ are independent and identically distributed
as
\[
    \bx_i \sim N(\mathbf{0},\,\bSig), \qquad i=1,\dots,n,
\]
where $\bSig\in\bbR^{p\times p}$ is a known, positive-definite,
block-diagonal covariance matrix.
The phenotype vector is
\[
    \y = \X\balpha + \beps,
\]
where $\balpha=(\alpha_1,\dots,\alpha_p)^\top$ is a \emph{fixed}
(non-random) vector of genetic effects and
$\beps=(\epsilon_1,\dots,\epsilon_n)^\top$ consists of independent
$N(0,\sigma_\epsilon^2)$ errors, independent of $\X$.

\medskip
\noindent\textbf{Key derived quantities.}

\begin{definition}[Marginal signal and noise]\label{def:beta}
Define
\[
    \beta_j \;=\; (\bSig\balpha)_j \;=\; \sum_{k=1}^{p}\Sigma_{jk}\,\alpha_k,
    \qquad j=1,\dots,p,
\]
and write $\bbeta = \bSig\balpha \in \bbR^p$.
Define also the marginal phenotypic variance
\[
    \sigma_y^2 \;=\; \balpha^\top\bSig\,\balpha + \sigma_\epsilon^2.
\]
\end{definition}

\begin{definition}[LD scores and estimator]\label{def:ld}
For $j=1,\dots,p$,
\begin{align}
    \ell_j &= \sum_{i=1}^{p}\Sigma_{ji}^2, &
    \widehat{a}_j &= n^{-1}\X_j^\top\y, &
    \widehat{w}_j &= \widehat{a}_j^{\,2}, \label{eq:defs}
\end{align}
where $\X_j$ denotes the $j$-th column of $\X$.
Set $\boldsymbol{\ell}=(\ell_1,\dots,\ell_p)^\top$,
$\widehat{\mathbf{w}}=(\widehat{w}_1,\dots,\widehat{w}_p)^\top$,
and
\[
    \mu_\ell = \frac{1}{p}\,\one_p^\top\boldsymbol{\ell},
    \qquad
    \widetilde{\boldsymbol{\ell}} = \boldsymbol{\ell} - \mu_\ell\,\one_p .
\]
The centered oracle LDSC estimator is
\begin{equation}\label{eq:tildesig}
    \widetilde\sigma_\alpha^2
    = \frac{\widetilde{\boldsymbol{\ell}}^{\,\top}
            (\widehat{\mathbf{w}} - \widehat\mu_w\,\one_p)}
           {\widetilde{\boldsymbol{\ell}}^{\,\top}\,
            \widetilde{\boldsymbol{\ell}}} .
\end{equation}
\end{definition}

\begin{remark}\label{rem:center}
Since
$\widetilde{\boldsymbol{\ell}}^{\,\top}\one_p
= \one_p^\top\boldsymbol{\ell} - p\,\mu_\ell = 0$,
the intercept--centering term $\widehat\mu_w\,\one_p$ drops out:
\begin{equation}\label{eq:reduced}
    \widetilde\sigma_\alpha^2
    = \frac{\widetilde{\boldsymbol{\ell}}^{\,\top}\widehat{\mathbf{w}}}
           {\|\widetilde{\boldsymbol{\ell}}\|^2}
    = \frac{\displaystyle\sum_{j=1}^{p}\widetilde\ell_j\,\widehat{a}_j^{\,2}}
           {\displaystyle\sum_{j=1}^{p}\widetilde\ell_j^{\,2}} .
\end{equation}
Because $\boldsymbol{\ell}$ depends only on the population covariance
$\bSig$, the denominator $\|\widetilde{\boldsymbol{\ell}}\|^2$ is
deterministic.
\end{remark}

%% ═══════════════════════════════════════════════════════════
\section{Preliminary Moment Calculations}
%% ═══════════════════════════════════════════════════════════

Throughout, we write $y_i = \bx_i^\top\balpha + \epsilon_i$ and
note that $(\bx_i,\epsilon_i)$, $i=1,\dots,n$, are mutually
independent.

\subsection{Representation as an i.i.d.\ average}
Define, for each sample $i$ and variant $j$,
\begin{equation}\label{eq:Tij}
    T_{ij} = X_{ij}\,y_i
           = X_{ij}\bigl(\bx_i^\top\balpha + \epsilon_i\bigr).
\end{equation}
Then $\widehat{a}_j = n^{-1}\sum_{i=1}^n T_{ij}$,
and for each fixed $j$ the random variables $T_{1j},\dots,T_{nj}$
are i.i.d.\ copies of
\begin{equation}\label{eq:Tj}
    T_j = X_j\,(\bx^\top\balpha + \epsilon),
\end{equation}
where $\bx\sim N(0,\bSig)$ and $\epsilon\sim N(0,\sigma_\epsilon^2)$
are independent.

\subsection{First and second moments of $T_j$}

\begin{lemma}\label{lem:EandCov}
For all $j,k\in\{1,\dots,p\}$,
\begin{align}
    \E[T_j] &= \beta_j,  \label{eq:ET}\\
    \gamma_{jk}
    \;:=\; \Cov(T_j,T_k)
    &= \Sigma_{jk}\,\sigma_y^2 + \beta_j\,\beta_k . \label{eq:gamma}
\end{align}
\end{lemma}

\begin{proof}
Since $\epsilon$ is independent of $\bx$ with mean zero,
\[
    \E[T_j]
    = \E[X_j\,\bx^\top\balpha]
    = \sum_m \alpha_m\,\E[X_j X_m]
    = \sum_m \alpha_m\,\Sigma_{jm}
    = \beta_j.
\]
For the second moment,
\begin{align*}
    \E[T_j\,T_k]
    &= \E\bigl[X_j\,X_k\,(\bx^\top\balpha+\epsilon)^2\bigr]
    = \E\bigl[X_j X_k\,(\bx^\top\balpha)^2\bigr]
      + \sigma_\epsilon^2\,\E[X_j X_k].
\end{align*}
By Isserlis' theorem (fourth moment of Gaussians),
\begin{align*}
    \E\bigl[X_j X_k\,(\bx^\top\balpha)^2\bigr]
    &= \sum_{m,l}\alpha_m\alpha_l\,\E[X_j X_k X_m X_l]\\
    &= \sum_{m,l}\alpha_m\alpha_l
       \bigl(\Sigma_{jk}\Sigma_{ml}
             +\Sigma_{jm}\Sigma_{kl}
             +\Sigma_{jl}\Sigma_{km}\bigr)\\
    &= \Sigma_{jk}\,(\balpha^\top\bSig\balpha)
       + 2\,\beta_j\,\beta_k.
\end{align*}
Hence
$\E[T_j T_k]
= \Sigma_{jk}(\balpha^\top\bSig\balpha + \sigma_\epsilon^2) + 2\beta_j\beta_k
= \Sigma_{jk}\sigma_y^2 + 2\beta_j\beta_k$,
and $\gamma_{jk} = \E[T_j T_k] - \beta_j\beta_k
= \Sigma_{jk}\sigma_y^2 + \beta_j\beta_k$.
\end{proof}

\subsection{Third cross-central moments of $T_j$}

\begin{lemma}\label{lem:third}
Define $U_j = T_j - \beta_j$. For all $j,k\in\{1,\dots,p\}$,
\begin{equation}\label{eq:rho}
    \rho_{jkk}
    \;:=\; \E[U_j\,U_k^2]
    = 4\,\sigma_y^2\,\Sigma_{jk}\,\beta_k
      + 2\,\sigma_y^2\,\Sigma_{kk}\,\beta_j
      + 2\,\beta_j\,\beta_k^2.
\end{equation}
\end{lemma}

\begin{proof}
Write $W_j = T_j = X_j\,y$.  Then
$\rho_{jkk} = \E[W_j W_k^2] - 2\beta_k\,\E[W_j W_k]
+ \beta_k^2\,\beta_j - \beta_j\,\E[W_k^2]
+ 2\beta_j\beta_k^2 - \beta_j\beta_k^2$.
Simplifying,
\begin{equation}\label{eq:rho_expand}
    \rho_{jkk}
    = \E[W_j W_k^2]
      - 2\beta_k\,\gamma_{jk}
      - \beta_j\,\gamma_{kk}
      - \beta_j\beta_k^2.
\end{equation}
Since $W_j W_k^2 = X_j X_k^2\,y^3$ and
$y^3 = (\bx^\top\balpha + \epsilon)^3$,
we expand using independence of $\epsilon$
and the fact that $\E[\epsilon]=\E[\epsilon^3]=0$:
\[
    \E[W_j W_k^2]
    = \E\bigl[X_j X_k^2\,(\bx^\top\balpha)^3\bigr]
      + 3\,\sigma_\epsilon^2\,\E[X_j X_k^2\,\bx^\top\balpha].
\]

\noindent\emph{Sixth--moment term.}
Setting $s = \bx^\top\balpha$, Isserlis' theorem applied to
$\E[X_j X_k X_k X_{m_1}X_{m_2}X_{m_3}]$
(summed against $\alpha_{m_1}\alpha_{m_2}\alpha_{m_3}$
over all 15 complete pairings of six Gaussian factors) yields,
after collecting terms,
\[
    \E[X_j X_k^2\,s^3]
    = 6\,\Sigma_{jk}\,\beta_k\,(\balpha^\top\bSig\balpha)
    + 3\,\Sigma_{kk}\,\beta_j\,(\balpha^\top\bSig\balpha)
    + 6\,\beta_j\,\beta_k^2.
\]

\noindent\emph{Fourth--moment term.}
Similarly,
\[
    \E[X_j X_k^2\,s]
    = \sum_m \alpha_m\,\E[X_j X_k^2 X_m]
    = 2\,\Sigma_{jk}\,\beta_k + \beta_j\,\Sigma_{kk}.
\]

Combining,
\[
    \E[W_j W_k^2]
    = 6\,\sigma_y^2\,\Sigma_{jk}\,\beta_k
    + 3\,\sigma_y^2\,\Sigma_{kk}\,\beta_j
    + 6\,\beta_j\,\beta_k^2.
\]
Substituting into \eqref{eq:rho_expand}
with $\gamma_{jk}=\Sigma_{jk}\sigma_y^2+\beta_j\beta_k$ and
$\gamma_{kk}=\Sigma_{kk}\sigma_y^2+\beta_k^2$,
one obtains \eqref{eq:rho} after straightforward cancellation.
\end{proof}

%% ═══════════════════════════════════════════════════════════
\section{Main Result}
%% ═══════════════════════════════════════════════════════════

\begin{theorem}[Variance of the oracle LDSC estimator]\label{thm:main}
Under the model and notation above, write
$\widetilde\ell_j = \ell_j - \mu_\ell$ and
$\|\widetilde{\boldsymbol{\ell}}\|^2
= \sum_{j=1}^p \widetilde\ell_j^{\,2}$.
Then
\begin{equation}\label{eq:var_exact}
\boxed{
    \Var(\widetilde\sigma_\alpha^2)
    = \frac{1}{\|\widetilde{\boldsymbol{\ell}}\|^4}
      \sum_{j=1}^{p}\sum_{k=1}^{p}
      \widetilde\ell_j\,\widetilde\ell_k\,
      \Cov\!\bigl(\widehat{a}_j^{\,2},\,\widehat{a}_k^{\,2}\bigr),
}
\end{equation}
where, with $\beta_j=(\bSig\balpha)_j$, \ 
$\gamma_{jk}=\Sigma_{jk}\sigma_y^2+\beta_j\beta_k$, \ 
$\sigma_y^2 = \balpha^\top\bSig\balpha+\sigma_\epsilon^2$, and
$\rho_{jkk}$ as in \eqref{eq:rho},
\begin{equation}\label{eq:cov_ajak}
    \Cov\!\bigl(\widehat{a}_j^{\,2},\,\widehat{a}_k^{\,2}\bigr)
    = \frac{4\,\beta_j\,\beta_k\,\gamma_{jk}}{n}
    + \frac{2\,\gamma_{jk}^2
            + 2\,\beta_j\,\rho_{jkk}
            + 2\,\beta_k\,\rho_{kjj}}{n^2}
    + \frac{\kappa_{jk}}{n^3},
\end{equation}
where
\begin{equation}\label{eq:kappa}
    \kappa_{jk}
    = \E[U_j^2\,U_k^2] - \gamma_{jj}\,\gamma_{kk} - 2\,\gamma_{jk}^2
\end{equation}
is the fourth-order joint cumulant, with $U_j = T_j - \beta_j$.
\end{theorem}

\begin{proof}
Because the denominator $\|\widetilde{\boldsymbol{\ell}}\|^2$ is
deterministic (Remark~\ref{rem:center}), equation \eqref{eq:var_exact}
is immediate from \eqref{eq:reduced}.
It remains to establish \eqref{eq:cov_ajak}.

\smallskip
\noindent\textbf{Step 1: Reduction to i.i.d.\ averages.}
Define $U_{ij} = T_{ij} - \beta_j$ and
$\bar U_j = n^{-1}\sum_{i=1}^n U_{ij}$, so that
$\widehat{a}_j = \beta_j + \bar U_j$.  Hence
\[
    \widehat{a}_j^{\,2} - \E[\widehat{a}_j^{\,2}]
    = 2\beta_j\,\bar U_j + (\bar U_j^2 - n^{-1}\gamma_{jj}).
\]
Writing the product
$(\widehat{a}_j^{\,2}-\E[\widehat{a}_j^{\,2}])
(\widehat{a}_k^{\,2}-\E[\widehat{a}_k^{\,2}])$
and taking expectations gives
\begin{equation}\label{eq:cov_decomp}
    \Cov(\widehat{a}_j^{\,2},\widehat{a}_k^{\,2})
    = 4\beta_j\beta_k\,\E[\bar U_j\bar U_k]
      + 2\beta_j\,\E[\bar U_j\bar U_k^2]
      + 2\beta_k\,\E[\bar U_j^2\bar U_k]
      + \bigl(\E[\bar U_j^2\bar U_k^2] - n^{-2}\gamma_{jj}\gamma_{kk}\bigr),
\end{equation}
using $\E[\bar U_j]=0$.

\smallskip
\noindent\textbf{Step 2: Moments of i.i.d.\ averages.}
Since $U_{1j},\dots,U_{nj}$ are i.i.d.\ centered random variables:

\begin{enumerate}
\item[(a)] $\E[\bar U_j\bar U_k] = n^{-1}\gamma_{jk}$.

\item[(b)] By expanding $\bar U_j\bar U_k^2 =
n^{-3}\sum_{i,i',i''}U_{ij}U_{i'k}U_{i''k}$ and using independence
across samples together with $\E[U_{ij}]=0$, the only nonvanishing
contributions arise when $i=i'=i''$; hence
$\E[\bar U_j\bar U_k^2] = n^{-2}\,\rho_{jkk}$.

\item[(c)] Expanding
$\bar U_j^2\bar U_k^2
= n^{-4}\sum_{i_1,i_2,i_3,i_4}U_{i_1 j}U_{i_2 j}U_{i_3 k}U_{i_4 k}$
and partitioning over the three types of
nonvanishing index configurations gives
\begin{align*}
    \E[\bar U_j^2\bar U_k^2]
    &= \frac{1}{n^3}\,\E[U_j^2 U_k^2]
     + \frac{n-1}{n^3}
       \bigl(\gamma_{jj}\gamma_{kk}+2\gamma_{jk}^2\bigr).
\end{align*}
The three surviving partitions of $(i_1,i_2,i_3,i_4)$ are:
all four equal ($n$ tuples, contributing $\E[U_j^2U_k^2]$);
$\{i_1=i_2\}\neq\{i_3=i_4\}$ ($n(n{-}1)$ tuples, contributing
$\gamma_{jj}\gamma_{kk}$);
and $\{i_1=i_3\}\neq\{i_2=i_4\}$ or
$\{i_1=i_4\}\neq\{i_2=i_3\}$ ($2n(n{-}1)$ tuples, each contributing
$\gamma_{jk}^2$).
Therefore
\[
    \E[\bar U_j^2\bar U_k^2] - n^{-2}\gamma_{jj}\gamma_{kk}
    = \frac{\kappa_{jk}}{n^3}
      + \frac{2\,\gamma_{jk}^2}{n^2},
\]
where $\kappa_{jk}=\E[U_j^2U_k^2]-\gamma_{jj}\gamma_{kk}-2\gamma_{jk}^2$.
\end{enumerate}

\noindent\textbf{Step 3: Assembly.}
Substituting (a)--(c) into \eqref{eq:cov_decomp}:
\[
    \Cov(\widehat{a}_j^{\,2},\widehat{a}_k^{\,2})
    = \frac{4\beta_j\beta_k\gamma_{jk}}{n}
    + \frac{2\beta_j\rho_{jkk}+2\beta_k\rho_{kjj}+2\gamma_{jk}^2}{n^2}
    + \frac{\kappa_{jk}}{n^3}.
\]
This is precisely \eqref{eq:cov_ajak}.
\end{proof}

%% ═══════════════════════════════════════════════════════════
\section{Explicit Expansion of the Variance}
%% ═══════════════════════════════════════════════════════════

We now substitute the closed-form expressions from
Lemmas~\ref{lem:EandCov}--\ref{lem:third} to write each
coefficient in~\eqref{eq:cov_ajak} purely in terms of the
allowed quantities $n,\bSig,\boldsymbol{\ell},\balpha,\sigma_\epsilon^2$.

\subsection{The $O(n^{-1})$ coefficient}
From \eqref{eq:gamma},
\begin{equation}\label{eq:order1coeff}
    4\,\beta_j\,\beta_k\,\gamma_{jk}
    = 4\,\sigma_y^2\,\Sigma_{jk}\,\beta_j\,\beta_k
      + 4\,\beta_j^2\,\beta_k^2.
\end{equation}

\subsection{The $O(n^{-2})$ coefficient}
From \eqref{eq:rho},
\begin{align}
    2\beta_j\,\rho_{jkk} + 2\beta_k\,\rho_{kjj}
    &= 16\,\sigma_y^2\,\Sigma_{jk}\,\beta_j\,\beta_k
       + 4\,\sigma_y^2\,(\Sigma_{jj}\beta_k^2 + \Sigma_{kk}\beta_j^2)
       + 8\,\beta_j^2\beta_k^2 . \notag
\end{align}
Adding $2\gamma_{jk}^2 = 2\sigma_y^4\Sigma_{jk}^2
+ 4\sigma_y^2\Sigma_{jk}\beta_j\beta_k + 2\beta_j^2\beta_k^2$,
\begin{equation}\label{eq:order2coeff}
    2\gamma_{jk}^2 + 2\beta_j\rho_{jkk} + 2\beta_k\rho_{kjj}
    = 2\,\sigma_y^4\,\Sigma_{jk}^2
      + 20\,\sigma_y^2\,\Sigma_{jk}\,\beta_j\beta_k
      + 4\,\sigma_y^2\bigl(\Sigma_{jj}\beta_k^2+\Sigma_{kk}\beta_j^2\bigr)
      + 10\,\beta_j^2\beta_k^2.
\end{equation}

\subsection{Combined formula}
\begin{proposition}[Expanded variance]\label{prop:expanded}
Define the Hadamard square
$\mathbf{D}=\bSig\circ\bSig\in\bbR^{p\times p}$, i.e.\
$D_{jk}=\Sigma_{jk}^2$,
and write $\bbeta = \bSig\balpha$, 
$\bbeta^{\circ 2} = (\beta_1^2,\dots,\beta_p^2)^\top$,
$\mathbf{d} = (\Sigma_{11},\dots,\Sigma_{pp})^\top$.
Then
\begin{align}
    \Var(\widetilde\sigma_\alpha^2)
    &= \frac{1}{\|\widetilde{\boldsymbol{\ell}}\|^4}
    \Biggl\{
    \frac{4\,\sigma_y^2}{n}\;
        (\widetilde{\boldsymbol{\ell}}\circ\bbeta)^\top
        \bSig\,
        (\widetilde{\boldsymbol{\ell}}\circ\bbeta)
    \;+\;
    \frac{4}{n}\;
        \bigl(\widetilde{\boldsymbol{\ell}}^{\,\top}
              \bbeta^{\circ 2}\bigr)^2
    \notag\\[6pt]
    &\qquad\quad
    +\;\frac{2\,\sigma_y^4}{n^2}\;
       \widetilde{\boldsymbol{\ell}}^{\,\top}
       \mathbf{D}\,
       \widetilde{\boldsymbol{\ell}}
    \;+\;\frac{20\,\sigma_y^2}{n^2}\;
       (\widetilde{\boldsymbol{\ell}}\circ\bbeta)^\top
       \bSig\,
       (\widetilde{\boldsymbol{\ell}}\circ\bbeta^{\circ 2}/\bbeta)
    \notag\\[4pt]
    &\qquad\quad
    +\;\frac{4\,\sigma_y^2}{n^2}\;
      \bigl(\widetilde{\boldsymbol{\ell}}^{\,\top}\bbeta^{\circ 2}\bigr)
      \bigl(\widetilde{\boldsymbol{\ell}}^{\,\top}\mathbf{d}\bigr)
    \;+\;\frac{4\,\sigma_y^2}{n^2}\;
      \bigl(\widetilde{\boldsymbol{\ell}}^{\,\top}\mathbf{d}\bigr)
      \bigl(\widetilde{\boldsymbol{\ell}}^{\,\top}\bbeta^{\circ 2}\bigr)
    \notag\\[4pt]
    &\qquad\quad
    +\;\frac{10}{n^2}\;
       \bigl(\widetilde{\boldsymbol{\ell}}^{\,\top}
             \bbeta^{\circ 2}\bigr)^2
    \;+\;\frac{1}{n^3}\sum_{j,k}
       \widetilde\ell_j\,\widetilde\ell_k\,\kappa_{jk}
    \Biggr\}.
    \label{eq:var_expanded_ugly}
\end{align}
Because the middle lines above contain a notational ambiguity when
$\beta_j=0$, we restate the formula more carefully using index
summations.  Substituting
\eqref{eq:order1coeff}--\eqref{eq:order2coeff} into
\eqref{eq:var_exact}--\eqref{eq:cov_ajak}:
\begin{align}
\boxed{
\begin{aligned}
    \Var(\widetilde\sigma_\alpha^2)
    &= \frac{1}{\|\widetilde{\boldsymbol{\ell}}\|^4}
    \Biggl\{
    \frac{4\sigma_y^2}{n}
      \sum_{j,k}\widetilde\ell_j\widetilde\ell_k\,
        \Sigma_{jk}\beta_j\beta_k
    \;+\;
    \frac{4}{n}
      \Bigl(\sum_j\widetilde\ell_j\beta_j^2\Bigr)^{\!2}
    \\[6pt]
    &\qquad
    +\frac{2\sigma_y^4}{n^2}
      \sum_{j,k}\widetilde\ell_j\widetilde\ell_k\,\Sigma_{jk}^2
    \;+\;
    \frac{20\sigma_y^2}{n^2}
      \sum_{j,k}\widetilde\ell_j\widetilde\ell_k\,
        \Sigma_{jk}\beta_j\beta_k
    \\[4pt]
    &\qquad
    +\frac{4\sigma_y^2}{n^2}
      \sum_{j,k}\widetilde\ell_j\widetilde\ell_k
        (\Sigma_{jj}\beta_k^2+\Sigma_{kk}\beta_j^2)
    \;+\;
    \frac{10}{n^2}
      \Bigl(\sum_j\widetilde\ell_j\beta_j^2\Bigr)^{\!2}
    \\[4pt]
    &\qquad
    +\;\frac{1}{n^3}
      \sum_{j,k}\widetilde\ell_j\widetilde\ell_k\,\kappa_{jk}
    \Biggr\}
\end{aligned}
}
\label{eq:var_final}
\end{align}
where $\sigma_y^2=\balpha^\top\bSig\balpha+\sigma_\epsilon^2$, \ 
$\beta_j=(\bSig\balpha)_j$, and the fourth-order
cumulant $\kappa_{jk}$ is defined in \eqref{eq:kappa}.
\end{proposition}

\medskip

Equivalently, using the compact notation
$\mathbf{D}=\bSig\circ\bSig$ and
$\mathbf{d}=\mathrm{diag}(\bSig)$,
the $O(n^{-2})$ factor
$\sum_{j,k}\widetilde\ell_j\widetilde\ell_k
(\Sigma_{jj}\beta_k^2+\Sigma_{kk}\beta_j^2)
= 2(\widetilde{\boldsymbol\ell}^\top\mathbf{d})
   (\widetilde{\boldsymbol\ell}^\top\bbeta^{\circ 2})$.

\begin{remark}[Interpretation of scale]
Each line of \eqref{eq:var_final} has a clear statistical origin:
\begin{itemize}
\item The $O(n^{-1})$ terms arise from the interaction between
      the signal $\bbeta=\bSig\balpha$ and both the genetic
      sampling noise and the environmental noise $\sigma_\epsilon^2$.
      They vanish when $\balpha=\mathbf{0}$.
\item The $O(n^{-2})$ terms arise from (i) the purely noise-driven
      variance $2\sigma_y^4\Sigma_{jk}^2/n^2$, (ii) signal--noise
      cross-terms involving third moments, and (iii) squared-signal
      contributions.
\item The $O(n^{-3})$ term $\kappa_{jk}$ captures residual
      non-Gaussianity of the quadratic forms in $\widehat{a}_j^2$;
      it is computable via Isserlis' theorem (eighth moments of
      the joint Gaussian $(\bx,\epsilon)$) but is typically
      negligible for moderately large $n$.
\end{itemize}
\end{remark}

\begin{remark}[Null case $\balpha=\mathbf{0}$]
When all genetic effects are zero, $\beta_j=0$ for all $j$ and
$\sigma_y^2=\sigma_\epsilon^2$.  The variance simplifies to
\[
    \Var(\widetilde\sigma_\alpha^2)\Big|_{\balpha=\mathbf{0}}
    = \frac{2\,\sigma_\epsilon^4}{n^2}\cdot
      \frac{\widetilde{\boldsymbol\ell}^\top\,\mathbf{D}\,
            \widetilde{\boldsymbol\ell}}
           {\|\widetilde{\boldsymbol\ell}\|^4}
    + O(n^{-3}),
\]
recovering the classical variance of a Haseman--Elston--type
regression under the null.
\end{remark}

\begin{remark}[Role of the block-diagonal structure]
The assumption that $\bSig$ is block-diagonal is used in two ways.
First, it ensures that $\ell_j=\sum_k\Sigma_{jk}^2$ and all
the sums involving $\Sigma_{jk}$ are effectively finite-range
(supported within blocks), so that all moments
$\E[U_j^{a}U_k^{b}]$ for $j,k$ in different blocks factor into
independent contributions.
Second, it guarantees that the fourth-order cumulant $\kappa_{jk}$
in \eqref{eq:kappa} vanishes whenever $j$ and $k$ lie in different
LD blocks, since in that case $U_j$ and $U_k$ are independent up
to the shared noise $\epsilon$, and the dependence through $\epsilon$
alone contributes only through lower-order terms already accounted
for. Thus the $O(n^{-3})$ sum in \eqref{eq:var_final}
is a within-block sum.
\end{remark}

\end{document}
